\documentclass{article}

% if you need to pass options to natbib, use, e.g.:
%     \PassOptionsToPackage{numbers, compress}{natbib}
% before loading neurips_2021

% ready for submission
\usepackage[preprint]{neurips_2021}

% to compile a preprint version, e.g., for submission to arXiv, add add the
% [preprint] option:
%     \usepackage[preprint]{neurips_2021}

% to compile a camera-ready version, add the [final] option, e.g.:
%     \usepackage[final]{neurips_2021}

% to avoid loading the natbib package, add option nonatbib:
%    \usepackage[nonatbib]{neurips_2021}

\usepackage[utf8]{inputenc} % allow utf-8 input
\usepackage[T1]{fontenc}    % use 8-bit T1 fonts
\usepackage[colorlinks=true]{hyperref}       % hyperlinks
\usepackage{url}            % simple URL typesetting
\usepackage{booktabs}       % professional-quality tables
\usepackage{amsfonts}       % blackboard math symbols
\usepackage{nicefrac}       % compact symbols for 1/2, etc.
\usepackage{microtype}      % microtypography
\usepackage{xcolor}         % colors
\usepackage{graphicx}       % Extended version of graphics
\usepackage{subcaption}

\newcommand{\todo}[1]{\textcolor{red}{\textbf{#1}}}

\graphicspath{ {./images/} }

\title{Is lottery fair?}

% The \author macro works with any number of authors. There are two commands
% used to separate the names and addresses of multiple authors: \And and \AND.
%
% Using \And between authors leaves it to LaTeX to determine where to break the
% lines. Using \AND forces a line break at that point. So, if LaTeX puts 3 of 4
% authors names on the first line, and the last on the second line, try using
% \AND instead of \And before the third author name.

\author{%
  Michal Horáček\\
  Matrikelnummer 6373382\\
  \texttt{michal.horacek@student.uni-tuebingen.de} \\
  \And
  Carson Zhang\\
  Matrikelnummer 6384481\\
  \texttt{carson.zhang@student.uni-tuebingen.de} \\
}

\begin{document}

\maketitle

\begin{abstract}
    We investigate the uniformity of various lotteries. 
    To detect potential tampering for both the German Lotto and New York Quick Draw lottery, we test the null hypothesis that the distribution of the minimum distance between
    lottery numbers, $d$, is as it would be in a uniformly sampled lottery.
    Furthermore, we evaluate lottery as a random number generator using the Diehard battery of tests.
\end{abstract}

\section{Introduction}
\label{sec:introduction}
Lottery/Hazard business is big many. Games such as lottery and dice are (probably?) one of the first application of randomness in human culture.
But for the same time organizers have cheated in lotteries, skewing the uniform distribution of drawn numbers expected by common sense for their
own monetary gain. We investigate the distribution of answers for the German Lotto lottery drawn from 1955 to the present day.

In the second part of this paper, we explore the "quality" of random numbers generated by lottery draws via the Diehard battery of tests~\cite{diehard}. These
tests were developed in the 1990's by George Marsaglia and are still in use today to test random number generators~\cite{modernRNG}.


\section{Dataset}
\label{sec:dataset}
Section \ref{sec:methods} investigates numbers drawn from the German lottery Lotto, from 1955 onwards. 
We have chosen this dataset because it is relevant to us and readers of this analysis as current residents of Germany, 
and it contains almost 70 years of data, which we believ is enough data to perform a meaningful analysis.

Furthermore, Drakakis, Taylor, and Rickard did not apply their tests to this data, so 
performing the test on the German Lotto data is a meaningful new result.

The dataset has been compiled by Johannes Friedrich, 
a software developer who has made the data publicly available via 
\href{https://github.com/JohannesFriedrich/LottoNumberArchive}{his Github repository}.

\todo{Perform verification of the correctness of the dataset by sampling from it and manually inspecting the data.}

\todo{Is there a dataset for a rigged lottery?} \href{https://notebook.community/JesseScott/Lotto649/lotto}{This} looks sketchy as hell.

However even 70 years of Lotto numbers is not sufficient to produce enough data for the diehard tests. Their author recommends
using 10 to 12 MB of random bits, which is substantially more than 28.4 KiB of Lotto numbers. Therefore we add winning numbers
from other lotteries to our dataset. By data volume, our data is dominated by the New York Quick Draw lottery, because it has
been drawn every 4 minutes for the last decade, adding up to to 25,663,100 numbers at our collection of the dataset.
Washington DC's Keno is quite similar, reaching 2,299,563 numbers over 3 years of existence. Our dataset is completed by
the "other" category, consisting of 1,917,890 winning numbers combined from lotteries drawn less often worldwide.

\begin{figure}
    \centering
    \includegraphics[width=\textwidth]{dataset_composition.pdf}
    \caption{Origin of numbers making up our dataset.}
    \label{fig:dataset}
\end{figure}

The "other" dataset comes from European or english-speaking countries. The reason is twofold - we have been looking
for them with English search queries and these nations are more likely to subscribe to ideas like open data and thus
offer such dataset as csv files. The European nationality stands for European transnational lotteries Eurojackpot and Euromillions.

\begin{figure}
    \centering
    \includegraphics[width=\textwidth]{other_composition.pdf}
    \caption{National composition of the "other" numbers.}
\end{figure}

In total, we have gathered winning numbers from 37 different lotteries.


\section{Distribution testing}
\label{sec:distribution_testing}
\begin{table}

\caption{Frequencies of $d$ statistic for German Lotto}
\centering
\begin{tabular}[t]{c|lllllll}
% \toprule
d & 1 & 2 & 3 & 4 & 5 & 6 & 7 and 8\\
\hline
Expected & 2310.596 & 1266.76 & 639.888 & 290.255 & 113.59 & 35.858 & 9.025\\
Observed & 2383 & 1247 & 615 & 273 & 105 & 35 & 8\\
% \bottomrule
\end{tabular}
\label{tab:frequencies_german}
\end{table}

For the German Lotto and NY Quick Draw, we conducted a hypothesis test where $H_0$ is true when lottery numbers are sampled with equal probability, and $H_A$ is true otherwise. We used Pearson's $\chi^2$ goodness-of-fit test because we can satisfy the assumptions and rules-of-thumb its usage, and it is the same test used in previous work, so our results can add on to theirs~\cite{Drakakis}. We used $\alpha = 0.05$ because it is popular and Drakakis et al. used the same $\alpha$.

\subsection{Minimum-distance statistic}

To create the frequency tables for the $\chi^2$ test (example in Table \ref{tab:frequencies_german}), we used the minimum distance statistic $d$ ~\cite{Drakakis} .
It is useful for detecting human tampering because humans do a
poor job of imitating its distribution under the null hypothesis. Boland and Pawitan ~\cite{trying_to_be_random} asked humans to simulate a lottery where $6$ numbers are sampled from $42$, and 
computed a $\chi^2$ goodness-of-fit statistic of $36.45$ when comparing the human-produced $d$ against the uniformly produced $d$ (assuming $p - 1$ degrees of freedom, in a hypothesis test this would have a p-value of $< 0.0000001$).Therefore, a low p-value could be consistent with human tampering.

We computed the expected distribution using the following formula~\cite{mindist_distribution}.

Let the lottery game be a sample $r_1,...,r_m$ of $m$ integers drawn from the integers $1,...,n$. 
(In the German Lotto, $n = 49$ and $m = 6$.)
Then, $d$ has the following definition and distribution.
\begin{equation}
    d = \min_{1 \leq i < j \leq m} | r_j - r_i |\\
\end{equation}
\begin{equation}
    P(d < k) = 1 - \frac{{n - (k - 1)(m - 1) \choose m}}{{n \choose m}}, \quad k = 1,..., \left\lfloor \frac{n - 1}{m - 1} \right\rfloor
\end{equation}

% \begin{table}

% \caption{Frequencies of $d$ statistic for German Lotto}
% \centering
% \begin{tabular}[t]{c|lllllll}
% % \toprule
% d & 1 & 2 & 3 & 4 & 5 & 6 & 7 and 8\\
% \hline
% Expected & 2310.596 & 1266.76 & 639.888 & 290.255 & 113.59 & 35.858 & 9.025\\
% Observed & 2383 & 1247 & 615 & 273 & 105 & 35 & 8\\
% % \bottomrule
% \end{tabular}
% \label{tab:frequencies}
% \end{table}

\subsection{The test statistic}

The $\chi^2$ goodness-of-fit test tests how well an expected discrete distribution fits
an observed discrete distribution. Under $H_0$, each 
lottery drawing is independent, which is a requirement for this test~\cite{openintrostats}. We combined the counts of adjacent categories when necessary to satisfy this rule of thumb for usage of the 
$\chi^2$ goodness-of-fit test: the expected frequency in each bin must be $\geq 5$~\cite{Drakakis}.

We define the $\chi^2$ statistic as follows: let $O_i$ be the observed frequency of category $i$, and $E_i$ be the expected frequency. Then, for our data,

\begin{equation}
    \chi^2 = \sum_{i=1}^p \frac{(O_i - E_i)^2}{E_i} \sim \chi^2 (p - 1)
\end{equation}

We use $p - 1$
because we estimated no parameters from our data~\cite{Drakakis}.

\subsection{Results}

For the German Lotto, we computed $\chi^2_6 = 5.357 $, with a p-value of $0.498$. For the NY Quick Draw, we computed $\chi^2_1 = 1.407$, with a p-value of $0.235$. These results are not significant at $\alpha = 0.05$ or 
any other common significance level. For both lotteries, the distribution of $d$ looks like a typical
distribution for a fair lottery. We do not suspect
human tampering from these results.

\subsection{Discussion}

We considered testing the goodness-of-fit of the frequency of the drawn numbers to a discrete uniform distribution, but the lottery numbers
are not independent: the numbers are sampled without replacement,
so they are multivariate hypergeometric under $H_0$. Thus, their frequencies have non-zero covariance~\cite{siegrist}. This also means sequences of lottery numbers will not be truly random.

Testing the distributions of other statistics, such as the sum of a lottery drawing, could provide additional insight into the fairness of the lottery, and into which numbers a player might choose. We provide an implementation of this test.

This analysis is appropriate for detecting human-generated lottery 
numbers, but it is unknown how well it will more sophisticated forms of unfairness, i.e. non-uniform sampling.

\section{Diehard tests}
\label{sec:diehard}
Fairness entails more than the question whether are lottery numbers from the expected distribution. For instance, the Kolmogorov-Smirnov
test used in the first part of this paper does not concern itself with the order the numbers are drawn. However if we saw a lottery whose
numbers were always drawn in a descending sequence, for example, we would become suspicious.

Thus a more comprehensive test is clearly required to establish a more detailed answer to our question. We approach this problem by reformulating
lottery as a process producing a stream of (supposedly) random numbers, which themselves are simply bit sequences. Under this formulation,
we can deploy standard statistical tests developed for testing random number generators: we have a file of one and zero bits and wish to investigate
if its bits are correlated, repeating with a period or other quantities undesirable for randomness.

A number of these test suites has been developed over time. Donald Knuth presented an initial set of empirical tests in the second volume of his computer science bible
The Art of Computer Programming in 1969. Many general cryptography textbooks such as \href{http://www.cacr.math.uwaterloo.ca/hac/}{Handbook of Applied Cryptography} 
or \href{http://www.wisdom.weizmann.ac.il/~oded/foc-vol1.html}{Foundations of Cryptography} contain multiple tests of their own. The american National Institute 
of Standards \& Technology has published a \href{https://nvlpubs.nist.gov/nistpubs/legacy/sp/nistspecialpublication800-22r1a.pdf}{guideline} discussing this matter too.

We decided to use the Diehard battery of tests, which was developed by the american statistician George Marsaglia in the nineties. While this package used
to be quite popular in its day, it has been superceded today by other suites, including its derivatives such as Dieharder or TestU01. \todo{Why did we pick this one?}

\begin{enumerate}
    \item About Diehard.
    \item The following part is divided into several section. In section 1, we discuss data gathering, processing and general creation of input files for Diehard. Special
    attention is given to the problem of attaining input file of sufficient size and the asymetric requirements of individual Diehard tests.
    \item Part 2 highlights result on Diehard suite, interprets them and compares them against commonly used PRNGs.
    \item Section 3 clarifies the shortcomings of above approach.
\end{enumerate}


\section{Conclusion}
\label{sec:conclusion}
We analyzed the uniformity of lottery results in two major ways. We performed a hypothesis test of the distribution of the minimum distance $d$, a statistic that is known to be poorly imitated by humans. Based on the results for the German Lotto and NY Quick Draw, there is no strong evidence for human tampering, but it is unknown how well this analysis may detect other forms of unfairness, and it is unknown how well this would inform a player's choices.

% The Diehard analysis was greatly limited by insufficient amount of data and in some cases unsuitable format in which the data was reported in (ascending-order draws). Out of the three datasets, we can only make a semi-conclusive statement about DC Keno. Here the evidence points towards the validity of the null hypothesis; that is, the randomness and therefore fairness of the lottery.

The Diehard analysis was greatly limited by having too little data, the without-replacement nature of lottery sampling, and some obviously non-random (sorted) data. Despite these concerns, we believe that the performance of DC Keno on the Diehard tests suggests the fairness of this lottery.


\newpage

\bibliography{sources.bib}
\bibliographystyle{unsrt}

\end{document}
